\chapter{Mở đầu}

\section{Bối cảnh nghiên cứu}

Trong các bài toán thao tác robot nhiều bước, thông tin quan sát cảm biến thu được từ môi trường tại một thời điểm tức thời thường không phản ánh đầy đủ toàn bộ diễn tiến tương tác đã xảy ra. Cụ thể, một vật thể có thể được nhấc lên và đặt trở lại vị trí ban đầu sau một chu kỳ thao tác, hoặc có thể bị che khuất hoàn toàn khỏi góc nhìn camera sau khi được đặt vào trong vật chứa. Trong những tình huống như vậy, các quan sát cảm biến bề mặt tại các thời điểm khác nhau có thể trở nên gần như đồng dạng, trong khi quyết định hành động tiếp theo lại đòi hỏi các phản ứng trái ngược: tiếp tục thao tác nếu mục tiêu chưa hoàn thành, hoặc duy trì trạng thái kết thúc khi chuỗi yêu cầu đã được thỏa mãn. Dưới góc nhìn quá trình quyết định Markov quan sát một phần POMDP \cite{pomdp1998}, hiện tượng suy biến quan sát này khiến chính sách điều khiển không thể phân giải lịch sử chỉ dựa trên các khung hình tức thời. Đồng thời, cấu trúc bài toán trong các bộ tiêu chuẩn như LIBERO-Mem \cite{liberomem2025} hay các khảo sát mở rộng gần đây của RoboMME \cite{robomme2026} cho thấy nhiều tác vụ thao tác có bản chất lặp lại chu kỳ hoặc hoàn nguyên trạng thái, từ đó đặt ra nhu cầu tất yếu về một cơ chế lưu giữ và tái hiện thông tin lịch sử phù hợp.

Tuy nhiên, việc tích hợp lịch sử vào các mô hình thị giác - ngôn ngữ - hành động VLA hiện nay gặp phải những rào cản kỹ thuật lớn, đặc biệt đối với các mô hình VLA nhỏ gọn. Các kiến trúc VLA gọn nhẹ vốn được thiết kế nhằm tối ưu hóa tần số suy luận điều khiển vòng kín và tiết kiệm tài nguyên tính toán trên robot, do đó chúng chịu những ràng buộc hết sức nghiêm ngặt về độ dài ngữ cảnh và dung lượng biểu diễn \cite{smolvla2025}. Việc mở rộng đầu vào bằng cách lưu giữ toàn bộ chuỗi video quan sát hoặc chuỗi vector trạng thái thô sẽ nhanh chóng làm bùng nổ chi phí tính toán và vượt quá giới hạn đầu vào cho phép của mô hình \cite{mem2026}. Ngược lại, nếu thực hiện nén lịch sử một cách tùy tiện, mô hình dễ đánh mất các mốc sự kiện quan trọng hoặc gặp phải hiện tượng suy diễn sai lệch. Do đó, bài toán đặt ra là: làm thế nào để chắt lọc những dữ kiện lịch sử có ý nghĩa với tác vụ, xác định chính xác thời điểm dữ kiện đó có đủ tin cậy để sử dụng, và chiếu thông tin đó vào không gian ngữ cảnh hạn hẹp của mô hình điều khiển một cách hiệu quả nhất.

Trong khảo sát ban đầu về các tác vụ có tính lặp hoặc phụ thuộc thứ tự, nghiên cứu nhận thấy tiến trình thực thi có thể được trừu tượng hóa một cách mạch lạc thông qua các mốc sự kiện chuyển đổi trạng thái, ví dụ như số chu kỳ tương tác đã hoàn thành hoặc sự thay đổi quan hệ không gian giữa các thực thể. Khi quan sát tức thời không đủ phân biệt tiến độ, một biểu diễn sự kiện có bằng chứng sẽ cung cấp tín hiệu tiến độ ngữ nghĩa rõ ràng, giúp mô hình VLA nhận biết đúng giai đoạn hiện tại. Cần lưu ý rằng trong các bài toán điều khiển liên tục, chính sách VLA thường dự đoán một chuỗi hành động action chunk liên tục và môi trường thực thi sẽ dừng khi đạt giới hạn thời gian hoặc thỏa mãn điều kiện kết thúc, chứ không thông qua lệnh dừng độc lập. Đây chính là tiền đề cốt lõi để khóa luận xây dựng một khung bộ nhớ sự kiện tổng quát, phục vụ cho các chính sách thao tác robot tay máy đơn dưới các ràng buộc khắt khe về tài nguyên.

\section{Khoảng trống nghiên cứu và định vị đề tài}

Các nghiên cứu gần đây về cơ chế bộ nhớ cho mô hình VLA đã khảo sát nhiều hướng biểu diễn khác nhau. MEM \cite{mem2026} kết hợp lịch sử video ngắn hạn với bản tóm tắt ngôn ngữ dài hạn theo hai thang thời gian. MemoryVLA \cite{memoryvla2026} duy trì một ngân hàng đặc trưng tri giác và nhận thức để truy xuất động theo ngữ cảnh. Embodied-SlotSSM \cite{liberomem2025} tập trung theo dõi trạng thái tiềm ẩn cho từng thực thể riêng biệt, trong khi khảo sát diện rộng RoboMME \cite{robomme2026} chỉ ra rằng hiệu quả của các dạng bộ nhớ phụ thuộc sâu sắc vào cấu trúc của từng họ tác vụ. Mặc dù các biểu diễn tiềm ẩn liên tục này có ưu thế về việc bảo toàn nhiều tín hiệu tri giác cho điều khiển, chúng thường đòi hỏi tài nguyên tính toán lớn và mang tính chất hộp đen, gây khó khăn cho việc kiểm chứng chính xác nguồn dữ liệu và thời điểm một dữ kiện lịch sử bắt đầu tác động đến quyết định hành động. Do đó, việc nghiên cứu một biểu diễn sự kiện tường minh, có thể truy xuất nguồn gốc bằng chứng và có thời điểm khả dụng xác định là một hướng tiếp cận cần thiết để nâng cao tính minh bạch và độ tin cậy của hệ thống.

Khoảng trống nghiên cứu thứ nhất liên quan đến cách biểu diễn tiến độ thực thi và độ tin cậy của bằng chứng. Khi huấn luyện trên dữ liệu mẫu tĩnh, mô hình VLA có thể vô tình học sự tương quan giả giữa thời gian trôi qua và tiến độ tác vụ. Khi triển khai vòng kín, nếu robot thao tác chậm hơn hoặc bị gián đoạn, sự tương quan thời gian này sẽ dẫn đến các "tiến độ ảo" \cite{clockvla2026}. Các nghiên cứu như EventSync \cite{eventsync2026} hay Clock VLA \cite{clockvla2026} đã bước đầu chỉ ra tầm quan trọng của việc neo mốc tiến trình vào các sự kiện vật lý thực tế. Tuy nhiên, các phương pháp hiện có vẫn thiếu một cơ chế tường minh để biểu thị mức độ đầy đủ của lịch sử khi quan sát cảm biến bị che khuất hoặc gián đoạn. Khóa luận hướng tới giải quyết khoảng trống này bằng cách xây dựng bộ nhớ sự kiện kèm theo bằng chứng và mốc thời gian khả dụng, đồng thời duy trì trạng thái chưa xác định khi bằng chứng không đủ, nhằm bảo đảm chính sách không hành động dựa trên những suy đoán thiếu căn cứ.

Khoảng trống nghiên cứu thứ hai nằm ở phương pháp luận đánh giá phân tách trong điều khiển vòng kín. Nhiều nghiên cứu hiện nay có xu hướng gộp chung việc đánh giá độ chính xác của cơ chế trích xuất dữ kiện với kết quả hoàn thành tác vụ cuối cùng của chính sách, hoặc đồng nhất việc giảm sai số dự đoán ngoại tuyến với sự cải thiện năng lực điều khiển thực tế. Thực tế cho thấy, một bộ trích xuất hoạt động tốt trên dữ liệu mẫu có thể thất bại khi chính sách tự sinh quỹ đạo vòng kín; ngược lại, việc giảm sai số dự đoán trên tập dữ liệu chưa chắc chuyển hóa thành tỷ lệ thành công cao do các lỗi tích lũy động học. Khóa luận định vị phương pháp luận đánh giá theo ba tầng độc lập: (1) độ trung thực và độ trễ của bộ ghi sự kiện từ dữ liệu cảm biến; (2) độ nhạy của chính sách đối với thông tin sự kiện ở cấp độ quyết định cục bộ; và (3) hiệu năng điều khiển vòng kín có đối chứng chặt chẽ. Sự phân định này đặc biệt cốt lõi khi muốn kiểm chứng giá trị thực của một biểu diễn bộ nhớ rút gọn trên nền mô hình VLA đã tiền huấn luyện.

\section{Mục tiêu và câu hỏi nghiên cứu}

Khóa luận tập trung nghiên cứu cách thức lưu giữ, xác thực và cung cấp lịch sử tương tác cho các chính sách VLA nhỏ gọn điều khiển robot tay máy đơn. Trong các tác vụ chuỗi dài mà quan sát hiện tại bị suy biến hoặc không đủ phân giải tiến độ, hệ thống cần cung cấp cho chính sách thông tin về các mốc sự kiện đã xảy ra cùng mức độ tin cậy của bằng chứng hỗ trợ, thông qua một giao diện chiếu ngôn ngữ phù hợp với giới hạn ngữ cảnh. Mục tiêu của nghiên cứu là xây dựng và đánh giá thực nghiệm một kiến trúc bộ nhớ sự kiện có bằng chứng, phân tích sự đánh đổi giữa lợi ích tiến độ và chi phí tính toán, đồng thời làm rõ khả năng thích ứng của các mô hình VLA nhỏ gọn.

Từ mục tiêu trên, nghiên cứu đặt ra ba câu hỏi khoa học cụ thể:

\textbf{RQ1. Giá trị của bộ nhớ dưới giới hạn ngữ cảnh.} Trong điều kiện cùng ngân sách ngữ cảnh và cùng bộ trọng số mô hình nền, việc cung cấp thông tin sự kiện tiến độ giúp cải thiện chất lượng quyết định hành động và tỷ lệ hoàn thành tác vụ ở mức độ nào so với chính sách chỉ dựa vào quan sát tức thời?

\textbf{RQ2. Độ tin cậy và khả năng xử lý lịch sử không đầy đủ.} Khi dòng quan sát cảm biến bị che khuất, gián đoạn hoặc xuất hiện tín hiệu mâu thuẫn, cơ chế kiểm định bằng chứng và biểu thị trạng thái chưa hoàn tất tác động thế nào đến hành vi của chính sách? Việc duy trì ranh giới bằng chứng giúp giảm thiểu các hành vi vượt mốc hoặc lặp chu kỳ sai lệch ra sao?

\textbf{RQ3. Khả năng duy trì kỹ năng cơ sở và tính mở rộng.} Việc thích ứng mô hình VLA với đầu vào điều kiện hóa từ bộ nhớ ảnh hưởng thế nào đến các kỹ năng thao tác cơ sở đã học trước đó? Cơ chế bộ nhớ đề xuất có khả năng mở rộng sang các dạng phụ thuộc lịch sử đa dạng và các kiến trúc mô hình khác nhau như thế nào?

Cần làm rõ rằng trong phạm vi triển khai thực nghiệm của khóa luận, một phần các câu hỏi này đã được kiểm chứng bằng số liệu đo đạc: đối sánh vòng kín F0/F1, khảo sát sai số dự đoán ngoại tuyến, can thiệp lời nhắc tại một truy vấn, và đánh giá duy trì kỹ năng trên LIBERO-10 với cả SmolVLA và X-VLA. Công thức thích ứng F0/F1 đã được áp dụng trên mô hình nền thứ hai X-VLA. Phép đánh giá này đo mức duy trì hiệu năng sau tinh chỉnh trên LIBERO-10 gốc bằng cách so sánh các bản đã tinh chỉnh với B0, tức checkpoint nền trước khi tinh chỉnh trong khóa luận. Các phần còn mở gồm quét nhiều mức ngân sách token có đánh giá tác động lên policy, đối chứng lịch sử video phi cấu trúc, và xác nhận đa seed trên roster đã khóa trước.

\section{Phạm vi và đóng góp}

Về mặt lý thuyết và phương pháp, khóa luận xây dựng một khung biểu diễn bộ nhớ sự kiện tổng quát cho các tác vụ thao tác có tính lặp, theo dõi danh tính và biến đổi quan hệ thực thể. Về thực nghiệm, đánh giá vòng kín có bộ nhớ tập trung vào T1 (một chu kỳ) và T3 (ba chu kỳ) của LIBERO-Mem \cite{liberomemSource}, chủ yếu trên SmolVLA. Khóa luận cũng đánh giá khả năng duy trì kỹ năng trên 10 tác vụ LIBERO-10 \cite{liberoSource} bằng SmolVLA và X-VLA, với bộ nhớ tắt ở các điều kiện X-VLA. Kiểm tra truy vấn theo sự kiện còn có cohort calibration trên nhóm trạng thái thứ hai, nhưng lượt chạy này dùng \texttt{writer\_override=true} sau khi preflight nghiêm ngặt phát hiện sai khác nguồn writer nên chỉ mang tính chẩn đoán. Các phép chẩn đoán trích xuất sự kiện trên T7-T10 cũng nằm trong phạm vi. Mọi kết quả tuân thủ giao thức riêng của nghiên cứu và không được xem là điểm LIBERO chính thức do khác biệt về cấu hình thực thi và điều kiện dừng. Đánh giá bộ nhớ trên X-VLA bằng rollout LIBERO-Mem và triển khai trên robot thực tế vẫn nằm ngoài phạm vi bằng chứng hiện tại.

Khóa luận xác lập ba đóng góp khoa học chính:

\textbf{Đóng góp thứ nhất là phát biểu bài toán tiến độ và sự nhập nhằng quan sát trong thao tác robot.} Khóa luận hình thức hóa hiện tượng suy biến quan sát trong các tác vụ chuỗi dài, nơi mà các trạng thái cảm biến đồng dạng có thể tương ứng với các yêu cầu hành vi đối lập, chẳng hạn như tiếp tục chu kỳ hoặc duy trì trạng thái đích. Nghiên cứu chỉ ra rằng thông tin lịch sử cần thiết không phải là toàn bộ diễn biến cảm biến thô mà là các mốc chuyển đổi trạng thái có bằng chứng hỗ trợ, đồng thời phân định rạch ròi giữa biến tiến độ ước lượng từ cảm biến với biến tiến độ đặc quyền dùng để đánh giá.

\textbf{Đóng góp thứ hai là đề xuất kiến trúc bộ nhớ sự kiện có kèm bằng chứng và giao diện chiếu ngôn ngữ.} Khóa luận xây dựng mô hình máy trạng thái hữu hạn cho phép ghi nhận các sự kiện chuyển đổi quan hệ thực thể từ luồng quan sát cảm biến và trạng thái robot, kèm theo mốc thời gian khả dụng và cờ chỉ báo khi lịch sử bị gián đoạn. Khóa luận phát biểu và chứng minh các bất biến về quy tắc đếm theo cặp sự kiện và tính không nhìn trước non-anticipative của nội dung bộ nhớ. Bộ nhớ sau đó được chiếu thành tóm tắt văn bản ngắn gọn, phù hợp với các mô hình VLA nhỏ gọn như SmolVLA \cite{smolvla2025}, kết hợp bộ dò mở Grounding DINO \cite{groundingdino2024} và kỹ thuật thích ứng tham số hiệu quả LoRA \cite{lora2021}.

\textbf{Đóng góp thứ ba là xây dựng quy trình thực nghiệm đa tầng từ dữ liệu ngoại tuyến đến điều khiển vòng kín.} Quy trình gồm đánh giá bộ ghi sự kiện trên quỹ đạo mẫu, đối sánh vòng kín, kiểm tra truy vấn theo sự kiện, can thiệp bỏ bộ nhớ tại một truy vấn, và đánh giá duy trì hiệu năng sau tinh chỉnh. Trong đối sánh vòng kín, F0 đạt 4/40 lượt thành công, F1 đạt 0/40; kiểm tra truy vấn theo sự kiện trên nhóm phát triển ghi nhận 1/24 lượt thành công ở mỗi arm. Nhóm hiệu chuẩn của phép kiểm tra truy vấn theo sự kiện sau override cũng ghi nhận 1/16 ở mỗi arm, nhưng chỉ là kết quả chẩn đoán vì preflight nghiêm ngặt bị chặn bởi sai khác hash nguồn writer. Các ca can thiệp cho thấy việc bỏ nội dung bộ nhớ tại một truy vấn làm thay đổi hành động và kết cục trong hai ca được chọn, chưa phải ước lượng hiệu quả tổng quát. Phép đánh giá này ghi nhận SmolVLA B0 đạt 37/50 so với F1-off 23/50; trên X-VLA, B0 đạt 47/50, F0 đạt 44/50 và F1-off đạt 46/50. Các cấu hình trong phép đánh giá này đều nhận chỉ dẫn gốc khi suy luận, kể cả F1-off, vốn giữ nguyên trọng số F1 đã tinh chỉnh với lời nhắc bộ nhớ. Toàn bộ mã nguồn, dữ liệu khóa và quy trình tái lập được tài liệu hóa trong hồ sơ kiểm định bằng chứng của đề tài.

Sau đây Chương 2 sẽ trình bày cơ sở về quyết định dưới quan sát một phần và đối chiếu các cách biểu diễn lịch sử trong VLA. Chương 3 xây dựng bộ nhớ sự kiện, phát biểu các tính chất và thiết kế phép đối chứng. Chương 4 phân tích độ phù hợp của sự kiện, kết quả ngoại tuyến và hành vi vòng kín theo các câu hỏi nghiên cứu. Chương 5 tổng hợp các kết luận trong phạm vi bằng chứng và xác định các thí nghiệm cần thiết để kiểm định giả thuyết còn mở.
